\chapter{Arnold relations}\label{app:arnold-relations}

\section{The affine relation}
\label{sec:arnold-historical}
\index{Arnold relations!detailed|textbf}

Arnold~\cite{Arnold69} computed the cohomology of
$C_n(\mathbb{A}^1)=\operatorname{Conf}_n(\mathbb{C})$. For
$1\leq i<j\leq n$ put
\[
\eta_{ij}=d\log(z_i-z_j)=\frac{dz_i-dz_j}{z_i-z_j}.
\]
We use the symmetric convention
$\eta_{ji}=\eta_{ij}$. Orlik--Solomon generators may instead be
chosen antisymmetric; the formulas below are written for the
differential forms themselves.

\begin{theorem}[Arnold relations \cite{Arnold69}; \ClaimStatusProvedElsewhere]
\label{thm:arnold-relations-appendix}
On $C_n(\mathbb{A}^1)$ the logarithmic forms $\eta_{ij}$ satisfy, for
distinct $i,j,k$,
\begin{equation}\label{eq:arnold-original}
\eta_{ij}\wedge\eta_{jk}
+\eta_{jk}\wedge\eta_{ki}
+\eta_{ki}\wedge\eta_{ij}=0.
\end{equation}
The cohomology algebra $H^*(C_n(\mathbb{A}^1);\mathbb{Q})$ is the
Orlik--Solomon algebra of the braid arrangement: it is generated by
the degree-one classes $\eta_{ij}$ subject to
$\eta_{ij}^2=0$ and the relations~\eqref{eq:arnold-original}.
\end{theorem}

\begin{proof}
Write $z_{ab}=z_a-z_b$ and $dz_{ab}=dz_a-dz_b$. The identity
\begin{equation}\label{eq:arnold-cyclic-appendix}
z_{ij}+z_{jk}+z_{ki}=0
\end{equation}
and its differential give
$dz_{ki}=-(dz_{ij}+dz_{jk})$. Hence
\[
dz_{jk}\wedge dz_{ki}=dz_{ij}\wedge dz_{jk},
\qquad
dz_{ki}\wedge dz_{ij}=dz_{ij}\wedge dz_{jk}.
\]
The three summands of~\eqref{eq:arnold-original} therefore have the
same numerator. With denominator $z_{ij}z_{jk}z_{ki}$ their coefficient
is $z_{ki}+z_{ij}+z_{jk}$, which is zero by
\eqref{eq:arnold-cyclic-appendix}. Arnold's computation, in the
Orlik--Solomon form of Brieskorn~\cite{Brieskorn73} and
Orlik--Solomon~\cite{OS80}, proves that these quadratic relations
generate the full ideal for the braid arrangement.
\end{proof}

\section{Affine, projective, and positive-genus relation packages}
\label{sec:arnold-genus-split-appendix}

The name ``Arnold relation'' refers here to the affine relation
\eqref{eq:arnold-original} on $C_n(\mathbb{A}^1)$ and on formal
genus-zero collision screens. The projective line and positive-genus
curves carry this relation on each tangent screen, together with extra
global relations.

\begin{proposition}[Configuration-space relation packages; \ClaimStatusProvedHere]
\label{prop:arnold-genus-split-appendix}
Let $X$ be a smooth complex curve.
\begin{enumerate}[label=\textup{(\roman*)}]
\item For $X=\mathbb{A}^1$, the forms $\eta_{ij}$ generate
$H^*(C_n(\mathbb{A}^1);\mathbb{Q})$, with only
$\eta_{ij}^2=0$ and Arnold's three-term relations.

\item For $X=\mathbb{P}^1$ and $n\geq3$,
\[
C_n(\mathbb{P}^1)
\cong
\mathrm{PSL}_2(\mathbb{C})\times
C_{n-3}(\mathbb{C}\setminus\{0,1\}).
\]
On the gauge-fixed factor, with moving points $z_4,\ldots,z_n$, put
\[
a_i=d\log z_i,\qquad
b_i=d\log(z_i-1),\qquad
c_{ij}=d\log(z_i-z_j).
\]
The degree-one Orlik--Solomon factor is governed by
\begin{align}
a_i\wedge b_i &=0,\label{eq:arnold-app-p1-rel-1}\\
a_i\wedge c_{ij}+c_{ij}\wedge a_j+a_j\wedge a_i&=0,\label{eq:arnold-app-p1-rel-2}\\
b_i\wedge c_{ij}+c_{ij}\wedge b_j+b_j\wedge b_i&=0,\label{eq:arnold-app-p1-rel-3}\\
c_{ij}\wedge c_{jk}+c_{jk}\wedge c_{ki}+c_{ki}\wedge c_{ij}&=0.\label{eq:arnold-app-p1-rel-4}
\end{align}
Equivalently, the pairwise projective classes
$\bar\eta_{ij}$ satisfy Arnold's three-term relations and
\[
\sum_{j\neq i}\bar\eta_{ij}=0.
\]
The exterior generator of $H^3(\mathrm{PSL}_2(\mathbb{C});\mathbb{Q})$
is separate from the pairwise logarithmic forms.

\item For $X=\Sigma_g$, $g\geq1$, the rational model contains the
surface classes from $H^*(\Sigma_g^n)$ and degree-one generators
\(\eta_{ij}\) with Totaro--Kriz--Bezrukavnikov relations
\[
\eta_{ij}(\pi_i^*\xi-\pi_j^*\xi)=0
\qquad
(\xi\in H^*(\Sigma_g;\mathbb{Q}))
\]
and differential
\begin{equation}\label{eq:arnold-app-totaro-differential}
d\eta_{ij}
=
\Delta_{ij}
=
\omega_i+\omega_j+
\sum_{r=1}^g
\bigl(\alpha_{r,i}\beta_{r,j}-\beta_{r,i}\alpha_{r,j}\bigr).
\end{equation}
Thus $\eta_{ij}$ are no longer closed cohomology classes by
themselves. The elliptic and prime-form propagators realize the same
defect analytically through period terms.
\end{enumerate}
\end{proposition}

\begin{proof}
Part~\textup{(i)} is Theorem~\ref{thm:arnold-relations-appendix}.
Part~\textup{(ii)} follows by the unique M\"obius transformation carrying
three ordered points to $(0,1,\infty)$, after which the displayed
relations are the Arnold relations for the triples
$(0,1,i)$, $(0,i,j)$, $(1,i,j)$, and $(i,j,k)$. Part~\textup{(iii)}
is Totaro's diagonal model for configuration spaces of smooth
projective varieties specialized to curves~\cite[Theorems~1 and~3]{Totaro96},
together with the equivalent rational models of
Kriz~\cite{Kriz94} and Bezrukavnikov~\cite{Bezrukavnikov94}.
\end{proof}

\section{The affine residue square}

\begin{proposition}[Affine screen residue cancellation; \ClaimStatusProvedHere]
\label{prop:operadic-equivalence-arnold}
On a genus-zero affine collision screen, Arnold's relation is equivalent
to the vanishing of the triple-collision scalar part of the simple-pole
residue square:
\begin{equation}\label{eq:arnold-residue-cyclic}
\operatorname{Res}_{D_{ij}}\operatorname{Res}_{D_{jk}}
+\operatorname{Res}_{D_{jk}}\operatorname{Res}_{D_{ki}}
+\operatorname{Res}_{D_{ki}}\operatorname{Res}_{D_{ij}}=0
\end{equation}
for every distinct triple $(i,j,k)$, with the Koszul signs determined
by the order of the logarithmic factors. For a KZ/Kohno connection,
flatness also requires the infinitesimal braid relations on the
coefficient operators.
\end{proposition}

\begin{proof}
The residue of $d\log t\wedge\alpha$ along $t=0$ is
\(\alpha|_{t=0}\). Applying this rule to
\[
\eta_{ij}\wedge\eta_{jk},\qquad
\eta_{jk}\wedge\eta_{ki},\qquad
\eta_{ki}\wedge\eta_{ij}
\]
identifies the three cyclic summands in~\eqref{eq:arnold-original}
with the three iterated residues in~\eqref{eq:arnold-residue-cyclic}.
The signs are exactly the signs needed to move the consumed
logarithmic factor to the left before taking the residue. Hence
Arnold's form identity and the scalar triple-residue cancellation are
the same identity written in the two languages. If the residues are
operator-valued, the same scalar calculation is multiplied by the
coefficient commutators; the remaining condition is precisely the
infinitesimal braid relation.
\end{proof}

\begin{theorem}[Affine bar square; \ClaimStatusProvedHere]
\label{thm:bar-d-squared-arnold}
Let $A$ be a chiral algebra whose local bar differential on a
genus-zero affine collision screen has simple-pole residue part
\[
d_{\mathrm{res}}=\sum_{i<j} d_{ij}.
\]
The triple-collision component of $d_{\mathrm{res}}^2$ vanishes by
Arnold's relation. The full identity $d_{\mathrm{bar}}^2=0$ is the sum
of this affine residue cancellation with the internal differential,
the Borcherds or Stasheff identities for the OPE operations, and the
projective or period terms of
Proposition~\ref{prop:arnold-genus-split-appendix} when the curve is
not an affine screen.
\end{theorem}

\begin{proof}
The square of the residue part separates into diagonal, disjoint, and
triple-index terms. The diagonal terms vanish because a logarithmic
simple pole has no second residue along the same divisor. Disjoint
pairs anticommute by the Koszul sign rule and independence of the two
normal directions. The shared-index terms are exactly the cyclic
triple sum~\eqref{eq:arnold-residue-cyclic}, hence vanish by
Proposition~\ref{prop:operadic-equivalence-arnold}. The remaining
summands of the bar differential are controlled by the corresponding
algebraic identities of the chiral algebra. On $\mathbb{P}^1$ and on
positive-genus curves the total differential includes the additional
relations listed in Proposition~\ref{prop:arnold-genus-split-appendix}.
\end{proof}

\begin{remark}[Higher coherences]
The three-term Arnold relation is the quadratic circuit relation of
the braid arrangement. Longer Orlik--Solomon circuit relations are
consequences of triangular circuits in the complete graph. In the
ordered associative bar complex, the boundary faces are consecutive
blocks; their codimension-two cancellations are Stasheff identities.
The consecutive-block data alone do not imply all pairwise Arnold
relations. The passage from ordered to symmetric data requires the
named compatibility supplied by the Arnold/Kohno flatness calculation
and, in the braided case, by a strong-unitary descent datum.
\end{remark}

\section{Orlik--Solomon presentation}
\label{sec:arnold-nine-term}

\begin{definition}[Orlik--Solomon algebra]\label{def:orlik-solomon-arnold}
\index{Orlik--Solomon!algebra}
For the braid arrangement
\(\mathcal{A}_{n-1}=\{H_{ij}:z_i=z_j\}\) in \(\mathbb{A}^n\), let
\[
\operatorname{OS}(\mathcal{A}_{n-1})
=
\bigwedge(e_{ij}:1\leq i<j\leq n)/I
\]
where \(I\) is generated by \(e_{ij}^2=0\) and
\[
e_{ij}\wedge e_{jk}
+e_{jk}\wedge e_{ki}
+e_{ki}\wedge e_{ij}=0
\]
for all distinct \(i,j,k\), in the symmetric-generator convention.
\end{definition}

\begin{lemma}[OS computes cohomology \cite{OS80}; \ClaimStatusProvedElsewhere]\label{lem:OS-cohomology-arnold}
The assignment \(e_{ij}\mapsto[\eta_{ij}]\) gives
\[
\operatorname{OS}(\mathcal{A}_{n-1})
\cong
H^*(C_n(\mathbb{A}^1);\mathbb{C}).
\]
\end{lemma}

\begin{proof}
The logarithmic forms generate \(H^1\) and satisfy the displayed
relations by Theorem~\ref{thm:arnold-relations-appendix}. The
Orlik--Solomon theorem~\cite{OS80} identifies the complement of any
complex hyperplane arrangement with the corresponding
Orlik--Solomon algebra; the braid arrangement gives the stated
configuration space.
\end{proof}

\begin{remark}[Triple count]\label{rem:nine-term-detailed-arnold}
For \(n=3,4,5\) the number of triangular Arnold relations is
\(\binom{n}{3}=1,4,10\). These are the degree-two circuit relations of
the braid arrangement. Higher-degree Orlik--Solomon relations come
from longer circuits and reduce to triangular relations because the
complete graph has chords.
\end{remark}

\begin{corollary}[Affine simple-pole bar differential; \ClaimStatusProvedHere]\label{cor:bar-d-squared-zero-arnold}
On an affine genus-zero screen, the simple-pole residue part of the
geometric bar differential squares to zero:
\[
d_{\mathrm{res}}^2=0.
\]
\end{corollary}

\begin{proof}
This is Theorem~\ref{thm:bar-d-squared-arnold} restricted to the
simple-pole residue summand on the affine screen.
\end{proof}

\subsection{Attribution}

Arnold~\cite{Arnold69} proved the braid-arrangement case. Brieskorn
extended the method to hyperplane arrangements~\cite{Brieskorn73}.
Orlik and Solomon gave the algebraic presentation~\cite{OS80}.
Beilinson--Drinfeld use the same local logarithmic geometry in the
chiral setting~\cite{BD04}. Totaro~\cite{Totaro96},
Kriz~\cite{Kriz94}, and Bezrukavnikov~\cite{Bezrukavnikov94} provide
the positive-genus cohomological replacement; Fay's identities
control the analytic prime-form corrections~\cite{Fay73}.

\section{Nilpotency on a fixed affine triple}
\label{sec:arnold-in-bar-nilpotency}

\begin{theorem}[Affine Arnold and triple-residue nilpotency; \ClaimStatusProvedHere]\label{thm:arnold-iff-nilpotent}
On a genus-zero affine collision screen, the following are equivalent
for a fixed distinct triple \(i,j,k\):
\begin{enumerate}
\item Arnold's relation
\[
\eta_{ij}\wedge\eta_{jk}
+\eta_{jk}\wedge\eta_{ki}
+\eta_{ki}\wedge\eta_{ij}=0.
\]
\item The cyclic simple-pole triple-residue obstruction vanishes:
\[
\operatorname{Res}_{D_{ij}}\operatorname{Res}_{D_{jk}}
+\operatorname{Res}_{D_{jk}}\operatorname{Res}_{D_{ki}}
+\operatorname{Res}_{D_{ki}}\operatorname{Res}_{D_{ij}}=0.
\]
\item The triple-index component of \(d_{\mathrm{res}}^2\) in the
affine simple-pole bar model vanishes.
\end{enumerate}
This equivalence is confined to the affine simple-pole screen. The
projective-line and positive-genus global differentials have the extra
terms of Proposition~\ref{prop:arnold-genus-split-appendix}.
\end{theorem}

\begin{proof}
The equivalence of \textup{(1)} and \textup{(2)} is
Proposition~\ref{prop:operadic-equivalence-arnold}. The residue square
on a fixed triple is the sum of the three cyclic two-step collisions,
with the same signs, so \textup{(2)} is exactly \textup{(3)}.
The last sentence is Proposition~\ref{prop:arnold-genus-split-appendix}.
\end{proof}

\subsection{Explicit residue calculation}

\begin{computation}[Residues of logarithmic forms]\label{comp:explicit-residues}
On \(C_3(\mathbb{A}^1)\) set
\[
\omega=\eta_{12}\wedge\eta_{23}
=d\log(z_1-z_2)\wedge d\log(z_2-z_3).
\]
Near \(D_{12}\) use coordinates
\[
u=z_2,\qquad \epsilon=z_1-z_2,\qquad v=z_3.
\]
Then
\[
\eta_{12}=\frac{d\epsilon}{\epsilon},
\qquad
\eta_{23}=d\log(u-v),
\]
and
\[
\operatorname{Res}_{D_{12}}(\omega)
=d\log(u-v)=\eta_{23}|_{D_{12}}.
\]
A second residue along \(D_{23}\) gives \(1\). The other cyclic
summands give the same scalar with the signs forced by the wedge order.
Their signed sum is zero precisely because
\[
\eta_{12}\wedge\eta_{23}
+\eta_{23}\wedge\eta_{31}
+\eta_{31}\wedge\eta_{12}=0.
\]
\end{computation}

\subsection{\texorpdfstring{Arnold relations for $n=4$}{Arnold relations for n=4}}

\begin{computation}[All Arnold relations for four affine points]\label{comp:arnold-n4}
For \(C_4(\mathbb{A}^1)\) the four triangular relations are
\begin{align*}
\eta_{12}\wedge\eta_{23}+\eta_{23}\wedge\eta_{31}
+\eta_{31}\wedge\eta_{12}&=0,\\
\eta_{12}\wedge\eta_{24}+\eta_{24}\wedge\eta_{41}
+\eta_{41}\wedge\eta_{12}&=0,\\
\eta_{13}\wedge\eta_{34}+\eta_{34}\wedge\eta_{41}
+\eta_{41}\wedge\eta_{13}&=0,\\
\eta_{23}\wedge\eta_{34}+\eta_{34}\wedge\eta_{42}
+\eta_{42}\wedge\eta_{23}&=0.
\end{align*}
They are independent in degree two. Since
\(\dim \bigwedge^2 H^1=\binom{6}{2}=15\), one gets
\[
\dim H^2(C_4(\mathbb{A}^1))=15-4=11,
\]
in agreement with
\((1+t)(1+2t)(1+3t)=1+6t+11t^2+6t^3\).
\end{computation}

\subsection{\texorpdfstring{The general affine pattern}{The general affine pattern}}

\begin{theorem}[Arnold relations for \texorpdfstring{$n$}{n} affine points \cite{Arnold69, OS80}; \ClaimStatusProvedElsewhere]\label{thm:arnold-general-n}
For \(n\) points on \(\mathbb{A}^1\), the degree-one cohomology has
basis \(\eta_{ij}\), \(1\leq i<j\leq n\), and
\[
\dim H^1(C_n(\mathbb{A}^1))=\binom{n}{2}.
\]
The triangular Arnold relations generate the degree-two ideal, and
\[
\dim H^2(C_n(\mathbb{A}^1))
=
\binom{\binom{n}{2}}{2}-\binom{n}{3}.
\]
The full Poincar\'e polynomial is
\[
P_t(C_n(\mathbb{A}^1))=\prod_{k=1}^{n-1}(1+kt).
\]
\end{theorem}

\begin{proof}
The Orlik--Solomon theorem for the braid arrangement gives the
presentation and the Poincar\'e polynomial. The degree-two formula is the
coefficient of \(t^2\): the exterior square of the degree-one span is
quotiented by the \(\binom{n}{3}\) independent triangular circuit
relations.
\end{proof}

\begin{remark}[OPE interpretation]\label{rem:ope-arnold}
On an affine genus-zero simple-pole screen, Arnold's relation is the
scalar form of the Jacobi identity for the chiral bracket. For full
OPEs the Borcherds identity controls all pole orders. For KZ/Kohno
connections, Arnold supplies the logarithmic-form part and the
infinitesimal braid relations supply the operator part. Consecutive
Stasheff associativity is the ordered partial-compactification shadow;
the all-pairwise Arnold/Kohno flatness condition is a separate
compatibility.
\end{remark}

\section{Global Fulton--MacPherson boundary relations}
\label{sec:global-boundary-relations}

\begin{theorem}[Configuration-space boundary relations; \ClaimStatusProvedHere]
\label{thm:config-boundary-relations}
Let \(X\) be a smooth algebraic curve and
\(\overline{C}_n(X)\) its Fulton--MacPherson compactification. The
boundary divisor is normal crossing. Its codimension-one components
are indexed by subsets \(S\subseteq\{1,\ldots,n\}\), \(|S|\geq2\), and
its nonempty codimension-two intersections are exactly nested or
disjoint pairs:
\[
S'\subset S,\qquad
S_1\cap S_2=\varnothing.
\]
Crossing pairs \(S\cap S'\neq\varnothing\), with neither subset
contained in the other, do not meet. With the orientation signs from
the ordered normal directions, the boundary operator satisfies
\[
\partial_{\mathrm{FM}}^2=0.
\]
\end{theorem}

\begin{proof}
Fulton--MacPherson~\cite{FM94} construct \(\overline{C}_n(X)\) by
iterated blowup along the diagonals. The resulting boundary is a
normal-crossing divisor. The intersection criterion is the standard
nested-set criterion for the wonderful compactification of the
diagonal arrangement. In a local normal-crossing chart with boundary
coordinates \(t_1,\ldots,t_r\), the two orders of restricting to
\(\{t_a=t_b=0\}\) differ by a transposition of normal directions, hence
contribute opposite signs to the cellular or logarithmic boundary
operator. Summing over all codimension-two strata gives
\(\partial_{\mathrm{FM}}^2=0\).
\end{proof}

\begin{corollary}[\texorpdfstring{$d_{\mathrm{res}}^2$}{d-res squared} and global corrections; \ClaimStatusProvedHere]\label{cor:dres-squared-global}
The affine pairwise residue summand satisfies \(d_{\mathrm{res}}^2=0\)
on every genus-zero tangent screen. The global FM differential on
\(\mathbb{P}^1\) includes the projective linear relations, and the
global differential on a positive-genus curve includes the
Totaro--period terms of
\eqref{eq:arnold-app-totaro-differential}. The total FM boundary
differential squares to zero in each case.
\end{corollary}

\begin{proof}
The tangent-screen assertion is
Corollary~\ref{cor:bar-d-squared-zero-arnold}. The total global
statement is Theorem~\ref{thm:config-boundary-relations} applied to
the relevant logarithmic form package. Proposition
\ref{prop:arnold-genus-split-appendix} lists the additional projective
and positive-genus summands, which confine the affine Arnold
presentation to its natural ambient.
\end{proof}

\begin{remark}[Nine local terms]\label{rem:nine-term-global}
For three affine points, \(d_{\mathrm{res}}^2\) expands into diagonal,
disjoint, and shared-index terms. The diagonal terms vanish, the
disjoint terms cancel by signs, and the shared-index terms form the
three Arnold summands. For \(n>3\), the same local calculation occurs
on each affine triple of a tangent screen. The global compactification
adds the nested-set boundary signs, and the curve type determines
whether projective or period terms also enter.
\end{remark}
